\begin{abstract}
Các agent dựa trên LLM (large language model) có thể sinh trọn vẹn một chương
trình giảng dạy (instructional program), nhưng chất lượng của chúng phụ thuộc vào
\emph{constructive alignment}---sự tương ứng giữa learning objective, activity và
assessment---vốn là thuộc tính sụp đổ đầu tiên khi các tổ chức chuyển từ model lớn
sang small language model (SLM) chi phí thấp để sinh nội dung đào tạo nội bộ. Suy
giảm alignment là một khiếm khuyết tích lũy và vô hình: chương trình vẫn "nhìn ổn"
theo từng phần và chỉ lộ lỗi khi triển khai. Các hướng LLM-agent hiện có cho
instructional systems design (ISD) chỉ \emph{khuyến khích} alignment qua prompt;
không hướng nào biểu diễn nó thành trạng thái mà hệ thống có thể kiểm tra và sửa,
nên chất lượng bị phó mặc cho năng lực thô của model. Chúng tôi đề xuất
\textbf{AlignmentGraph-ISD}, một graph-guided multi-agent harness hiện thực hoá
constructive alignment thành một \emph{alignment graph} tường minh và kiểm tra được,
rồi tổ chức việc sinh xoay quanh nó: các agent chuyên biệt xây dựng graph theo một workflow cố định với
phạm vi ghi bị giới hạn, nên alignment hình thành \emph{by construction}, và mỗi
agent chạy pedagogical verification cùng repair như một công cụ ngay trong lượt của
mình. Chúng tôi đánh giá harness trên ISD-Agent-Bench, so với sáu agent tiêu biểu,
dưới cả một model lớn lẫn một SLM 9B, và thực hiện một ablation để tách riêng đóng
góp của từng cơ chế. Trên benchmark, AlignmentGraph-ISD đạt điểm tổng cao nhất
trong bảy agent; ablation cho thấy multi-agent decomposition trên graph là cơ chế
đóng góp chính ($+3{,}9$ điểm tổng, $p < 10^{-12}$), và---quan trọng---đóng
góp này gần như gấp đôi khi model nhỏ đi (từ $+3{,}9$ điểm tổng trên LLM lên $+7{,}9$
trên SLM, $p < 10^{-12}$), trong khi điểm của chính harness gần như bất biến theo
scale còn các baseline suy giảm. Pedagogical verification/repair không làm thay đổi
điểm do judge chấm trên cả hai model dù giảm được số alignment defect đo được. Các
kết quả này cho thấy: biểu diễn alignment thành trạng thái graph hạng nhất, và phân
rã việc sinh xoay quanh nó, chính là điều cho phép một agent harness duy trì chất
lượng thiết kế giảng dạy độc lập với model bên dưới.
\end{abstract}
